\resumeSubHeadingListStart
    \resumeSubheading
    {MLOps Engineer, MLOps}{Oct 2021 -- Jun 2022}
    {}{중구, 서울}

    
    \resumeProjectHeading
    {\textbf{미술시장 분석 PoC}}{Apr 2022 -- Jun 2022}
    \resumeItemListStart
        \resumeItem{HuggingFace NLP의 개체명 인식으로 아티스트 CV 데이터를 정량 분석하고 정제}
        \resumeItem{Word Histogram으로 기사 분석}
        \resumeItem{감성 분석으로 기사 분석}
        \resumeItem{\textbf{NLP로 텍스트 분석}}
    \resumeItemListEnd
    \hypertarget{mlops_back}{
    \resumeProjectHeading
    {\textbf{인스타그램 관계 데이터 분석}}{Feb 2022 -- Apr 2022}
    }
    \resumeItemListStart
        \resumeItem{Node2Vec · K-means · PCA로 관계 분석}
        \resumeItem{NextJS · Flask · CytoscapeJS로 관계를 방향 그래프로 보여주는 시각화 페이지 개발}
        \resumeItem{분석 자료를 토대로 Series A 투자 유치}
        \resumeItem{\textbf{모델 추론에 넣을 형태로 데이터를 정제하고 피처를 추출해 추론 결과를 분석}}
        \resumeItem{\hyperlink{mlops}{\underline{Project Detail >>}}}
    \resumeItemListEnd

    \hypertarget{instagram_scraper_back}{
    \resumeProjectHeading
    {\textbf{인스타그램 관계 정보 스크래퍼 개발}}{Oct 2021 -- Jun 2023}
    }
    \resumeItemListStart
        \resumeItem{데이터를 끊김 없이 모으기 위해 여러 방법을 구상하고 구현}
        \resumeItem{Selenium으로 동적 페이지 스크래퍼 개발}
        \resumeItem{requests로 Public API를 호출하는 스크래퍼 개발}
        \resumeItem{Instagram private API를 다루는 오픈소스 Instagrapi로 스크래퍼 개발}
        \resumeItem{Python으로 된 Instagrapi를 Swift로 옮겨 모바일에서 도는 스크래퍼 개발}
        \resumeItem{
            소셜 관계 정보를 모으다 오픈소스에 기여: \href{https://github.com/subzeroid/instagrapi/graphs/contributors}{\underline{Instagrapi}}
        }
        \resumeItem{\textbf{막힌 문제를 될 때까지 다른 방법으로 두드림}}
        \resumeItem{\hyperlink{instagram_scraper}{\underline{Project Detail >>}}}
    \resumeItemListEnd

    \resumeProjectHeading
    {\textbf{미술 아티스트 정보 스크래퍼 개발}}{Oct 2021 -- Nov 2021}
    \resumeItemListStart
        \resumeItem{requests로 정적 페이지 스크래퍼 개발}
        \resumeItem{boto3와 멀티스레딩으로 EC2 인스턴스를 여러 대 띄워 고속 스크래핑 환경 구성}
        \resumeItem{아티스트 정보 60만 건 · CV 정보 300만 건 수집}
        \resumeItem{\textbf{이후 분석의 바탕이 된 데이터 확보}}
    \resumeItemListEnd

    \hypertarget{postgresql_back}{
    \resumeProjectHeading
    {\textbf{PostgreSQL 기반 수집 데이터 관리}}{}
    }
    \resumeItemListStart
        \resumeItem{스키마를 나눠 원본 데이터와 정제 데이터를 구분해 보관}
        \resumeItem{Python과 psycopg2로 데이터 정제}
        \resumeItem{\textbf{초기 형태의 데이터 파이프라인 구축}}
        \resumeItem{\hyperlink{postgresql}{\underline{Project Detail >>}}}
    \resumeItemListEnd
\resumeSubHeadingListEnd
